\documentclass[12pt,a4paper,openany,oneside]{book}

% --- Paquetes ---
\usepackage[utf8]{inputenc}
\usepackage[spanish, es-noshorthands]{babel}
\usepackage{amsmath, amsfonts, amssymb}
\usepackage{graphicx}
\usepackage{geometry}
\usepackage{hyperref}
\usepackage{booktabs}
\usepackage{fancyhdr}
\usepackage{listings}
\usepackage{xcolor}
\usepackage{tikz}
\usepackage{pgfplots}
\usepackage{tikz}
\usetikzlibrary{matrix}
\usepackage{caption} % para \captionof

\usetikzlibrary{arrows.meta, positioning, shapes.geometric, calc}
\pgfplotsset{compat=1.18}

\geometry{margin=2.5cm}

% --- Colores y estilos para código Python ---
\definecolor{codegreen}{rgb}{0,0.6,0}
\definecolor{codegray}{rgb}{0.5,0.5,0.5}
\definecolor{codepurple}{rgb}{0.58,0,0.82}
\definecolor{backcolour}{rgb}{0.95,0.95,0.92}

\lstset{
    language=Python,
    backgroundcolor=\color{backcolour},   
    commentstyle=\color{codegreen},
    keywordstyle=\color{blue},
    numberstyle=\tiny\color{codegray},
    stringstyle=\color{codepurple},
    basicstyle=\ttfamily\small,
    breaklines=true,
    frame=single,
    showstringspaces=false
}

% --- Estilo de cabecera y pie ---
\pagestyle{fancy}
\fancyhf{}
\renewcommand{\headrulewidth}{0.4pt}
\renewcommand{\footrulewidth}{0.4pt}
\fancyhead[L]{\small Carlos César Sánchez Coronel}
\fancyhead[R]{\small Sesión 4: Sistemas de Ecuaciones y Regresión}
\fancyfoot[L]{\small Carlos César Sánchez Coronel}
\fancyfoot[C]{\thepage}

% --- Portada ---
\title{
    \textbf{Sesión 4} \\ 
    \Huge \textbf{SISTEMAS DE ECUACIONES Y LA BASE DE LA REGRESIÓN} \\
    \large De la eliminación gaussiana a los mínimos cuadrados \\
}
\author{
    \small Por \\ \vspace{0.2cm}
    \Large \textsc{Carlos César Sánchez Coronel}
}
\date{2026}

\begin{document}

\maketitle
\tableofcontents
\addcontentsline{toc}{chapter}{Índice general}

% ------------------------------------------------------------------
% CAPÍTULO 4: SESIÓN 4 - SISTEMAS DE ECUACIONES Y REGRESIÓN
% ------------------------------------------------------------------
\chapter{Sistemas de Ecuaciones y la Base de la Regresión}

En esta sesión abordaremos uno de los problemas fundamentales del álgebra lineal: la resolución de sistemas de ecuaciones lineales. Este conocimiento es la base de la regresión lineal, el primer modelo predictivo que todo estudiante de machine learning debe dominar. Veremos cómo representar un sistema en forma matricial, cómo resolverlo mediante eliminación gaussiana y qué hacer cuando el sistema tiene más ecuaciones que incógnitas (sistemas sobredeterminados), donde la solución por mínimos cuadrados se convierte en la herramienta ideal.

\section{Representación matricial de sistemas lineales}

Un sistema de $m$ ecuaciones lineales con $n$ incógnitas $x_1, x_2, \dots, x_n$ se escribe como:

\begin{align*}
a_{11}x_1 + a_{12}x_2 + \cdots + a_{1n}x_n &= b_1 \\
a_{21}x_1 + a_{22}x_2 + \cdots + a_{2n}x_n &= b_2 \\
&\vdots \\
a_{m1}x_1 + a_{m2}x_2 + \cdots + a_{mn}x_n &= b_m
\end{align*}

Este sistema se puede expresar en forma matricial como:

\begin{equation*}
\mathbf{A}\mathbf{x} = \mathbf{b}
\end{equation*}

donde:
\begin{itemize}
    \item $\mathbf{A} \in \mathbb{R}^{m \times n}$ es la matriz de coeficientes (cada fila corresponde a una ecuación),
    \item $\mathbf{x} \in \mathbb{R}^{n}$ es el vector de incógnitas,
    \item $\mathbf{b} \in \mathbb{R}^{m}$ es el vector de términos independientes.
\end{itemize}

\begin{center}
\begin{tikzpicture}
\matrix[matrix of math nodes, left delimiter={[}, right delimiter={]}] (A) at (0,0) {
a_{11} & a_{12} & \cdots & a_{1n} \\
a_{21} & a_{22} & \cdots & a_{2n} \\
\vdots & \vdots & \ddots & \vdots \\
a_{m1} & a_{m2} & \cdots & a_{mn} \\
};
\matrix[matrix of math nodes, left delimiter={[}, right delimiter={]}] (x) at (3,0.5) {
x_1 \\ x_2 \\ \vdots \\ x_n \\
};
\node[scale=2] at (1.8,0) {$\cdot$};
\node[scale=2] at (4.5,0) {$=$};
\matrix[matrix of math nodes, left delimiter={[}, right delimiter={]}] (b) at (6,0) {
b_1 \\ b_2 \\ \vdots \\ b_m \\
};
\end{tikzpicture}
\captionof{figure}{Representación matricial de un sistema lineal.}
\end{center}

\subsection{Interpretación geométrica}
Cada ecuación representa un hiperplano en $\mathbb{R}^n$. La solución del sistema es la intersección de todos esos hiperplanos. Para $n=2$, cada ecuación es una recta en el plano; la solución es el punto donde se cortan.

\begin{center}
\begin{tikzpicture}
\matrix[matrix of math nodes, left delimiter={[}, right delimiter={]}] (A) at (0,0) {
a_{11} & a_{12} & \cdots & a_{1n} \\
a_{21} & a_{22} & \cdots & a_{2n} \\
\vdots & \vdots & \ddots & \vdots \\
a_{m1} & a_{m2} & \cdots & a_{mn} \\
};

\matrix[matrix of math nodes, left delimiter={[}, right delimiter={]}] (x) at (3,0.5) {
x_1 \\ x_2 \\ \vdots \\ x_n \\
};

\node[scale=2] at (1.8,0) {$\cdot$};
\node[scale=2] at (4.5,0) {$=$};

\matrix[matrix of math nodes, left delimiter={[}, right delimiter={]}] (b) at (6,0) {
b_1 \\ b_2 \\ \vdots \\ b_m \\
};
\end{tikzpicture}
\captionof{figure}{Representación matricial de un sistema lineal.}
\end{center}

\section{Resolución de sistemas cuadrados ($m=n$)}

Cuando el número de ecuaciones es igual al número de incógnitas y la matriz $\mathbf{A}$ es invertible, existe una única solución dada por $\mathbf{x} = \mathbf{A}^{-1}\mathbf{b}$. Sin embargo, calcular la inversa es costoso e inestable numéricamente. En la práctica se utilizan métodos directos como la eliminación gaussiana o factorizaciones.

\subsection{Eliminación gaussiana}
El método consiste en transformar el sistema original en uno equivalente triangular superior mediante operaciones elementales (intercambiar filas, multiplicar una fila por un escalar, sumar un múltiplo de una fila a otra). Luego se resuelve por sustitución regresiva.

\textbf{Ejemplo:} Resolver el sistema
\begin{align*}
2x + y - z &= 8 \\
-3x - y + 2z &= -11 \\
-2x + y + 2z &= -3
\end{align*}

Escribimos la matriz aumentada $[\mathbf{A}|\mathbf{b}]$:
\[
\left[\begin{array}{ccc|c}
2 & 1 & -1 & 8 \\
-3 & -1 & 2 & -11 \\
-2 & 1 & 2 & -3
\end{array}\right]
\]

Aplicamos eliminación:
\begin{enumerate}
    \item Usamos el pivote $a_{11}=2$. Hacemos ceros debajo:
        \begin{align*}
            \text{Fila2} &\leftarrow \text{Fila2} + \frac{3}{2}\text{Fila1} \\
            \text{Fila3} &\leftarrow \text{Fila3} + \text{Fila1}
        \end{align*}
    Obtenemos:
    \[
    \left[\begin{array}{ccc|c}
    2 & 1 & -1 & 8 \\
    0 & 0.5 & 0.5 & 1 \\
    0 & 2 & 1 & 5
    \end{array}\right]
    \]
    \item Intercambiamos filas 2 y 3 para tener pivote en (2,2):
    \[
    \left[\begin{array}{ccc|c}
    2 & 1 & -1 & 8 \\
    0 & 2 & 1 & 5 \\
    0 & 0.5 & 0.5 & 1
    \end{array}\right]
    \]
    \item Hacemos cero en (3,2): Fila3 $\leftarrow$ Fila3 - 0.25$\times$Fila2
    \[
    \left[\begin{array}{ccc|c}
    2 & 1 & -1 & 8 \\
    0 & 2 & 1 & 5 \\
    0 & 0 & 0.25 & -0.25
    \end{array}\right]
    \]
\end{enumerate}

Ahora resolvemos de abajo arriba:
\begin{align*}
0.25 z &= -0.25 \Rightarrow z = -1 \\
2y + 1(-1) &= 5 \Rightarrow 2y = 6 \Rightarrow y = 3 \\
2x + 3 - (-1) &= 8 \Rightarrow 2x = 4 \Rightarrow x = 2
\end{align*}

Solución: $\mathbf{x} = (2,3,-1)$.

\subsection{Factorización LU}
La eliminación gaussiana puede interpretarse como una factorización de $\mathbf{A}$ en el producto de una matriz triangular inferior $\mathbf{L}$ (con unos en la diagonal) y una matriz triangular superior $\mathbf{U}$: $\mathbf{A} = \mathbf{L}\mathbf{U}$. Una vez obtenida la factorización, resolver $\mathbf{A}\mathbf{x} = \mathbf{b}$ equivale a resolver primero $\mathbf{L}\mathbf{y} = \mathbf{b}$ (sustitución progresiva) y luego $\mathbf{U}\mathbf{x} = \mathbf{y}$ (sustitución regresiva). Es más eficiente cuando se deben resolver múltiples sistemas con la misma matriz $\mathbf{A}$.

\subsection{Resolución con NumPy}
Para sistemas cuadrados con solución única, NumPy proporciona \texttt{np.linalg.solve} que utiliza factorización LU con pivotamiento.

\begin{lstlisting}[language=Python]
import numpy as np

A = np.array([[2, 1, -1],
              [-3, -1, 2],
              [-2, 1, 2]], dtype=float)
b = np.array([8, -11, -3])

x = np.linalg.solve(A, b)
print("Solución:", x)  # [ 2.  3. -1.]

# Verificar
print("A @ x - b:", np.dot(A, x) - b)  # casi cero
\end{lstlisting}

\section{Sistemas sobredeterminados ($m > n$) y mínimos cuadrados}

En la mayoría de aplicaciones de machine learning tenemos más ecuaciones (datos) que incógnitas (parámetros). Por ejemplo, en regresión lineal con $n$ características y $m$ muestras ($m \gg n$), queremos encontrar un vector $\mathbf{x}$ que satisfaga aproximadamente $\mathbf{A}\mathbf{x} \approx \mathbf{b}$. El sistema no tiene solución exacta, así que buscamos la mejor solución en el sentido de minimizar la norma del error.

\subsection{Problema de mínimos cuadrados}
Dado $\mathbf{A} \in \mathbb{R}^{m \times n}$ con $m > n$ y $\mathbf{b} \in \mathbb{R}^m$, queremos encontrar $\mathbf{x}$ que minimice:
\begin{equation*}
\|\mathbf{A}\mathbf{x} - \mathbf{b}\|_2^2 = \sum_{i=1}^{m} \left( \sum_{j=1}^{n} A_{ij}x_j - b_i \right)^2
\end{equation*}

La solución se obtiene resolviendo las \textbf{ecuaciones normales}:
\begin{equation*}
\mathbf{A}^\top \mathbf{A} \mathbf{x} = \mathbf{A}^\top \mathbf{b}
\end{equation*}
Si $\mathbf{A}^\top \mathbf{A}$ es invertible (lo que ocurre si las columnas de $\mathbf{A}$ son linealmente independientes), entonces:
\begin{equation*}
\mathbf{x} = (\mathbf{A}^\top \mathbf{A})^{-1} \mathbf{A}^\top \mathbf{b}
\end{equation*}

Esta solución es la que proporciona la regresión lineal ordinaria (OLS).

\subsection{Interpretación geométrica}
La solución de mínimos cuadrados proyecta el vector $\mathbf{b}$ sobre el espacio columna de $\mathbf{A}$. El residuo $\mathbf{r} = \mathbf{b} - \mathbf{A}\mathbf{x}$ es ortogonal a dicho espacio.

\begin{center}
\begin{tikzpicture}
\draw[->] (0,0) -- (4,0) node[right] {col 1};
\draw[->] (0,0) -- (0,4) node[above] {col 2};
\draw[fill=blue!20] (0,0) -- (3,1) -- (2,3) -- cycle;
\node at (1.5,1.5) {$\mathrm{Col}(\mathbf{A})$};
\draw[->, thick, red] (0,0) -- (3,3) node[above right] {$\mathbf{b}$};
\draw[->, thick, blue] (0,0) -- (2.5,1.5) node[right] {$\mathbf{A}\mathbf{x}$};
\draw[dashed, thick] (3,3) -- (2.5,1.5) node[midway, right] {$\mathbf{r}$};
\end{tikzpicture}
\captionof{figure}{Proyección ortogonal de $\mathbf{b}$ sobre el espacio columna.}
\end{center}

\subsection{Ejemplo: Ajuste de una recta por mínimos cuadrados}
Queremos ajustar una recta $y = \beta_0 + \beta_1 x$ a los puntos $(x_i, y_i)$ con $i=1,\dots,m$. El sistema es:
\[
\begin{pmatrix}
1 & x_1 \\
1 & x_2 \\
\vdots & \vdots \\
1 & x_m
\end{pmatrix}
\begin{pmatrix}
\beta_0 \\ \beta_1
\end{pmatrix}
\approx
\begin{pmatrix}
y_1 \\ y_2 \\ \vdots \\ y_m
\end{pmatrix}
\]
Resolvemos por mínimos cuadrados.

\begin{lstlisting}[language=Python]
import numpy as np
import matplotlib.pyplot as plt

# Datos de ejemplo
x = np.array([0, 1, 2, 3, 4, 5])
y = np.array([1, 3, 2, 5, 4, 6])

# Construir matriz A (columna de unos para el intercepto)
A = np.vstack([np.ones_like(x), x]).T  # shape (6,2)

# Resolver ecuaciones normales
ATA = A.T @ A
ATb = A.T @ y
beta = np.linalg.solve(ATA, ATb)
print(f"Intercepto: {beta[0]:.4f}, Pendiente: {beta[1]:.4f}")

# Graficar
plt.scatter(x, y, label='Datos')
x_line = np.linspace(0, 5, 100)
y_line = beta[0] + beta[1] * x_line
plt.plot(x_line, y_line, 'r-', label='Ajuste mínimos cuadrados')
plt.xlabel('x')
plt.ylabel('y')
plt.legend()
plt.grid(True)
plt.show()
\end{lstlisting}

\subsection{Alternativa: descomposición QR}
Resolver las ecuaciones normales puede ser numéricamente inestable si $\mathbf{A}^\top \mathbf{A}$ está mal condicionada. Una alternativa más robusta es utilizar la descomposición QR: $\mathbf{A} = \mathbf{Q}\mathbf{R}$, donde $\mathbf{Q}$ tiene columnas ortonormales y $\mathbf{R}$ es triangular superior. Entonces la solución es $\mathbf{x} = \mathbf{R}^{-1}\mathbf{Q}^\top \mathbf{b}$.

NumPy implementa mínimos cuadrados con \texttt{np.linalg.lstsq}, que por defecto usa la descomposición QR.

\begin{lstlisting}[language=Python]
beta_lstsq, residuals, rank, s = np.linalg.lstsq(A, y, rcond=None)
print("Con lstsq:", beta_lstsq)
\end{lstlisting}

\subsection{Conexión con scikit-learn}
La regresión lineal de scikit-learn también resuelve mínimos cuadrados (por defecto usando la descomposición QR o SVD). El resultado es el mismo.

\begin{lstlisting}[language=Python]
from sklearn.linear_model import LinearRegression

model = LinearRegression()
model.fit(x.reshape(-1,1), y)
print("Intercepto sklearn:", model.intercept_)
print("Pendiente sklearn:", model.coef_[0])
\end{lstlisting}

\section{Aplicaciones en IA y cursos futuros}

\subsection{Machine Learning}
La regresión lineal es el modelo más simple y fundacional. Aparece como caso particular de redes neuronales (una neurona sin activación) y como componente en modelos más complejos. La solución por mínimos cuadrados es también la base de la regresión ridge cuando se añade regularización L2.

\subsection{Robótica}
En robótica, los sistemas de ecuaciones lineales aparecen en cinemática directa e inversa, en el control de manipuladores y en la estimación de estados (filtros de Kalman).

\subsection{Optimización}
Los problemas de mínimos cuadrados son un caso particular de optimización convexa. Muchos algoritmos de machine learning (como SVM, regresión logística) se formulan como problemas de optimización que en ciertos casos se reducen a sistemas lineales.

\subsection{Procesamiento de señales}
En filtrado adaptativo, se utilizan mínimos cuadrados recursivos (RLS) para estimar parámetros de filtros en tiempo real.

\section{Ejercicio integrador: Comparación de métodos}

A continuación, generamos un conjunto de datos con ruido y comparamos la solución de mínimos cuadrados obtenida por diferentes métodos: ecuaciones normales, descomposición QR y scikit-learn.

\begin{lstlisting}[language=Python]
import numpy as np
from sklearn.linear_model import LinearRegression
import time

# Generar datos sintéticos: y = 2 + 3x + ruido
np.random.seed(42)
m = 1000
x = np.random.randn(m, 1)
X = np.hstack([np.ones((m,1)), x])  # matriz con columna de unos
beta_true = np.array([2, 3])
y = X @ beta_true + np.random.randn(m,1)*0.5

# Método 1: Ecuaciones normales (solución analítica)
start = time.time()
beta_norm = np.linalg.inv(X.T @ X) @ X.T @ y
t_norm = time.time() - start

# Método 2: lstsq (QR)
start = time.time()
beta_lstsq, _, _, _ = np.linalg.lstsq(X, y, rcond=None)
t_lstsq = time.time() - start

# Método 3: sklearn
start = time.time()
model = LinearRegression(fit_intercept=False)  # ya tenemos columna de unos
model.fit(X, y)
beta_sk = model.coef_.flatten()
t_sk = time.time() - start

print("Verdaderos:     ", beta_true)
print("Ecuaciones normales:", beta_norm.flatten())
print("lstsq (QR):       ", beta_lstsq.flatten())
print("sklearn:          ", beta_sk)
print("Tiempos: Norm={:.6f}s, lstsq={:.6f}s, sklearn={:.6f}s".format(t_norm, t_lstsq, t_sk))
\end{lstlisting}

Observa que los tres métodos dan resultados prácticamente idénticos. Las diferencias de tiempo pueden ser apreciables para matrices grandes.

\section{Resumen}

En esta sesión hemos aprendido:
\begin{itemize}
    \item La representación matricial de sistemas lineales $\mathbf{A}\mathbf{x} = \mathbf{b}$.
    \item Métodos de resolución para sistemas cuadrados: eliminación gaussiana y factorización LU.
    \item El problema de mínimos cuadrados para sistemas sobredeterminados, su solución mediante ecuaciones normales y descomposición QR.
    \item La conexión directa con la regresión lineal en machine learning.
    \item Implementaciones eficientes en NumPy y scikit-learn.
\end{itemize}
Estos conceptos son fundamentales para entender modelos predictivos, optimización y muchas otras áreas de la inteligencia artificial.

\end{document}